% ============================================================================
% dip_proposal_v2.tex -- Reframed project proposal.
% Build: xelatex dip_proposal_v2.tex
% ============================================================================
\documentclass[12pt,a4paper]{article}
\usepackage[margin=1.8cm]{geometry}
\usepackage{amsmath,amssymb}
\usepackage{graphicx}
\usepackage{booktabs}
\usepackage{float}
\usepackage{array}
\usepackage{tabularx}
\usepackage{xcolor}
\usepackage{enumitem}
\usepackage{hyperref}
\usepackage{caption}
\usepackage{fontspec}
\setlength{\textfloatsep}{8pt plus 2pt minus 2pt}
\setlength{\intextsep}{6pt plus 2pt minus 2pt}
\setlength{\abovecaptionskip}{4pt}
\setlength{\belowcaptionskip}{2pt}
\setmainfont{Latin Modern Roman}

\sloppy

\title{\textbf{Digital Image Processing --- Team Project Proposal}\\[2pt]
\large Restoring Visibility Through Haze, Water, Glass and Plastic:
A Comparative Study of Classical Methods, with the Dark Channel Prior
(DCP) as the Baseline}
\author{Noel Jacob \quad Shaik Areesh \quad Sri Harshit K}
\date{CS7.404 Digital Image Processing, Monsoon 2026 \\ \today}

\begin{document}
\maketitle

\begin{abstract}
\noindent
Camera visibility is degraded either by the medium between the scene and
the sensor (atmospheric haze, turbid water) or by a protective or
intervening element in front of the lens (fogged or dirty glass covers,
translucent or aged plastic). This project asks which classical,
non-learning image-processing methods restore visibility best in each
case, why they succeed or fail, and how far a synthetically generated
degradation can stand in for a real capture. The Dark Channel Prior
(DCP) is used throughout as the reference method and as the lens for
explaining results; no machine learning is used in any restoration or
degradation step.

A real, self-captured dataset across the four media is compared against
a calibrated synthetic-degradation model (single-layer and stacked
multi-layer), with DCP plus a small set of classical and medium-aware
alternatives run on both. The deliverable list is kept lean -- ten
items -- by grouping closely related work under a single heading. Three
tracks (real data, synthetic data, restoration) run in parallel and
converge on a single comparison and analysis spine, mapped against the
course's own Report~1 / Report~2 / presentation dates.
\end{abstract}

\tableofcontents
\bigskip

% ============================================================================
\section{Problem Statement and Objectives}
% ============================================================================

Camera visibility is degraded either by the medium between the scene and
the sensor (atmospheric haze, turbid water) or by a protective or
intervening element in front of the lens (fogged or dirty glass covers,
translucent or aged plastic). These degradations differ physically, but
all reduce contrast and visibility in similar ways. This project asks
which classical, non-learning image-processing methods restore
visibility best in each case, why they succeed or fail, and how far a
synthetically generated degradation can stand in for a real capture. The
Dark Channel Prior (DCP) is used throughout as the reference method and
as the lens for explaining results. No machine learning is used in any
restoration or degradation step.

\begin{itemize}[nosep]
  \item Implement the Dark Channel Prior from scratch as the baseline
        restoration method.
  \item Build a real, self-captured dataset with clear references across
        four media (haze, water, glass, plastic) at three severity
        levels.
  \item Compare DCP against classical and medium-aware alternatives on
        every medium.
  \item Build a synthetic degradation model (single-layer and stacked
        multi-layer) and validate it against real captures.
  \item Explain successes and failures using dark-channel statistics,
        transmission behaviour and colour-cast analysis.
  \item Recommend a simple, medium-aware restoration pipeline for each
        medium, backed by the results.
\end{itemize}

% ============================================================================
\section{Data: Real, Synthetic and Existing Image Sets}
\label{sec:data}
% ============================================================================

Real captures are the primary evidence. Matched synthetic images
(degrading our own clear references with the forward model, then
comparing against the real degraded capture of the same scene) validate
that model. Bulk synthetic data, generated from public depth datasets,
drives severity sweeps and stacked-layer experiments that would be
impractical to capture by hand.

\begin{table}[H]
\centering
\normalsize
\begin{tabular}{@{}lp{3.6cm}p{3.9cm}p{4.3cm}@{}}
\toprule
Medium & Real (own captures) & Synthetic (approximation model) & Existing online image sets \\
\midrule
Haze & Outdoor haze/fog shots plus a clear reference on a clear day &
  $t=\exp(-\beta d)$; $A$ near white &
  O-HAZE / I-HAZE (real haze, paired GT); RESIDE SOTS (paired); RESIDE
  ITS / NYU Depth v2 (depth, for synthesis) \\
Water & Aquarium/tub with a reference shot; milk drops for turbidity
  levels &
  Per-channel $t_c=\exp(-\beta_c d)$, $\beta_R>\beta_G>\beta_B$;
  blue-green $A$ &
  UIEB for real underwater images (pseudo-reference only; rely on
  no-reference metrics) \\
Glass & Fog by steam/spray, dirt by dust/smudges, on a clear pane &
  Near-uniform low-frequency veil (blurred noise), droplet patterns,
  alpha-blended smudge masks &
  None fit -- captures and synthesis carry this medium \\
Plastic & Clear, translucent and aged plastic sheets over a scene &
  Uniform low $t$, tint gains, Gaussian blur &
  None fit -- captures and synthesis carry this medium \\
Stacked & Water/haze scene viewed through a fogged or dirty pane &
  Two-layer model: $I_2 = J t_1 t_2 + A_1 t_2(1-t_1) + A_2(1-t_2)$ &
  --- \\
\bottomrule
\end{tabular}
\caption{Data tiers by medium. Real captures are the primary evidence;
matched and bulk synthetic data extend coverage where capture is
impractical.}
\label{tab:data}
\end{table}

\textbf{Caveat.} Synthetic haze generated with the same model DCP
inverts favours DCP by construction, so the synthetic-real gap is
reported explicitly rather than hidden. Synthesis parameters are
calibrated on a few scenes and evaluated on held-out ones. Depth for our
own scenes comes from measured distances, a planar ramp, or 2--3
hand-defined layers -- no learned depth model is used.

% ============================================================================
\section{Pipeline: Hand-Implemented vs.\ Library}
\label{sec:pipeline}
% ============================================================================

The rule: anything under study or central to the project's claims is
implemented by hand; generic, uncontested primitives use a library.
Library use for baselines should be confirmed with the instructor.

\begin{table}[H]
\centering
\small
\begin{tabular}{@{}p{6.3cm}p{3.6cm}p{4.3cm}@{}}
\toprule
Stage & By hand & Library \\
\midrule
Dark channel (patch-min over channels) & Yes -- pipeline logic is ours &
  \texttt{scipy.ndimage} or \texttt{cv2.erode} (cross-check) \\
Atmospheric light, raw transmission, recovery ($t_0$ floor) & Yes & --- \\
Guided filter refinement & Optional manual version &
  \texttt{cv2.ximgproc.guidedFilter} (opencv-contrib) \\
Soft matting (stretch goal) & --- & \texttt{pymatting} \\
Medium-aware variants: UDCP/red channel, per-channel or local $A$,
  large-patch/near-global $t$ & Yes -- modifications of our DCP & --- \\
Gray-world / simple white balance & Manual gains &
  \texttt{cv2.xphoto} as cross-check \\
Homomorphic filtering (glass) & Yes -- log, FFT, high-pass, exp &
  NumPy FFT \\
Wiener deblur (plastic) & PSF estimation manual &
  \texttt{skimage.restoration.wiener} \\
CLAHE, histogram equalization, gamma & --- & OpenCV \\
Retinex (MSR, stretch) & Yes -- Gaussian-blur based & --- \\
Median / bilateral cleanup & --- & OpenCV \\
Forward degradation models (all media, stacked) & Yes &
  NumPy / SciPy for noise \& blur helpers \\
PSNR, SSIM, CIEDE2000 & --- & \texttt{skimage.metrics}, \texttt{skimage.color} \\
UCIQE, entropy, contrast, dark-channel histograms, model-fit residual,
  ranking correlation & Yes -- short formulas, part of the analysis &
  \texttt{scipy.stats} for Spearman/Kendall \\
Plots and tables & --- & Matplotlib, pandas \\
\bottomrule
\end{tabular}
\caption{Implementation ownership by pipeline stage.}
\label{tab:pipeline}
\end{table}

\textbf{Generalized skeleton with per-medium presets:} preprocess (none
/ white balance) $\rightarrow$ medium-light estimate (global /
per-channel / local) $\rightarrow$ dark-channel variant $\rightarrow$
transmission (raw, then guided filter) $\rightarrow$ recover
$\rightarrow$ postprocess (CLAHE or contrast stretch). A config file
selects the modules per medium. CAP is excluded because its
coefficients were fit by supervised learning.

\begin{table}[H]
\centering
\normalsize
\begin{tabular}{@{}ll@{}}
\toprule
Medium & Variants tested \\
\midrule
Haze & DCP, DCP + guided filter, Retinex, CLAHE \\
Water & White balance then DCP, UDCP, CLAHE \\
Glass & Large-patch/near-global $t$, local $A$, homomorphic filtering, CLAHE \\
Plastic & White balance then DCP, Wiener deblur, Retinex, CLAHE \\
\bottomrule
\end{tabular}
\caption{Medium-aware variant selection.}
\label{tab:variants}
\end{table}

% ============================================================================
\section{Deliverables: Order and Dependencies}
\label{sec:deliverables}
% ============================================================================

The deliverable list is kept lean -- ten items instead of a longer
breakdown -- by grouping closely related work (for example, capture
protocol with the dataset it produces, and the forward-model spec with
the code and calibration that implement it) under a single deliverable.
Three tracks run in parallel -- real data (D1), synthetic data
(D2--D3) and restoration (D4--D6) -- and converge at D7. From D7 onward
the pipeline is strictly serial through to the final deliverables.

\begin{figure}[H]
\centering
\includegraphics[width=0.86\textwidth]{figs/deliverable_dependency.png}
\caption{Deliverable dependency graph. Blue = real-data track, amber =
synthetic-data track, green = restoration track, grey = the serial
spine from experiment runs to final deliverables.}
\label{fig:spine}
\end{figure}

\textbf{Runs in parallel:} D1, D2$\rightarrow$D3, and D4$\rightarrow$D5
alongside D6, are three independent tracks until they all feed D7.
Restoration code (track B) is tested on O-HAZE / RESIDE images while it
waits for real and synthetic data to land.

\textbf{Strictly serial:} D1 (freeze) and D3 $\rightarrow$ D7
$\rightarrow$ D8 $\rightarrow$ D9 $\rightarrow$ D10.

\begin{table}[H]
\centering
\normalsize
\begin{tabular}{@{}p{5.5cm}p{1.6cm}p{7.3cm}@{}}
\toprule
Deliverable & Status & Why \\
\midrule
D4 DCP core \& refinement, D5 medium-aware variants & \textbf{Green} &
  Well-documented algorithm, small code surface \\
D6 Metrics \& baselines & \textbf{Green} & Mostly library calls \\
D2--D3 Synthesis pipeline \& synthetic data & \textbf{Amber} &
  Matching real captures needs tuning \\
D8 Real-vs-synthetic comparison & \textbf{Amber} &
  Depends on both data tracks landing on time \\
D9 Failure \& stacked-layer analysis & \textbf{Amber} & Needs extra captures \\
Soft matting \& Wiener tuning (within D5) & \textbf{Amber} &
  Slow / fiddly -- first to cut if behind \\
D1 Real dataset & \textbf{Red} &
  Capture logistics, lighting and tripod consistency -- main schedule risk \\
\bottomrule
\end{tabular}
\caption{Risk assessment (RAG status) by deliverable.}
\label{tab:rag}
\end{table}

% ============================================================================
\section{Work Division}
\label{sec:work}
% ============================================================================

\begin{table}[H]
\centering
\normalsize
\begin{tabular}{@{}p{3.0cm}p{7.7cm}p{3.8cm}@{}}
\toprule
Member & Owns & Deliverables \\
\midrule
A -- Data \& captures & Capture protocol, all real scenes and severity
  control, depth setups, dataset organisation, downloads (O-HAZE, UIEB,
  ITS/NYU) & D1, later stacked real shots, demo images \\
B -- Restoration \& evaluation & DCP from scratch, refinement,
  medium-aware variants, baselines, metrics module, experiment runner &
  D4--D7, D9 \\
C -- Synthesis \& comparison & Forward models, calibration, matched \&
  bulk synthetic data, real-vs-synthetic statistics and ranking
  agreement & D2, D3, D8, stacked-layer experiment \\
All & Understand and be able to explain the DCP core; report sections;
  rehearsal & D10 \\
\bottomrule
\end{tabular}
\caption{Work division across the three team members.}
\label{tab:work}
\end{table}

A few hours per week per person is realistic for this scope, provided
the real dataset is frozen early (around Oct 14) so the busiest
coursework weeks are not competing with capture logistics.

% ============================================================================
\section{Timeline Against Course Dates}
\label{sec:timeline}
% ============================================================================

\begin{table}[H]
\centering
\normalsize
\begin{tabular}{@{}p{2.8cm}p{11.2cm}@{}}
\toprule
Dates & Milestone \\
\midrule
Sep 28 -- 30 & Report 1: framing, pipeline diagram, capture protocol,
  one or two DCP outputs \\
Oct 1 -- 14 & Captures (haze \& glass, then water \& plastic), DCP +
  refinement, synthesis code. Data frozen around Oct 14 \\
Oct 15 -- 26 & Variants, baselines, metrics, calibration, first
  real-vs-synthetic comparison. Report 2 on Oct 26 \\
Oct 27 -- Nov 8 & Ranking agreement, failure \& stacked-layer analysis,
  severity sweeps, recommended pipelines, README \\
Nov 9 -- 12 & Rehearsal. Presentation on Nov 12 (Nov 9 is a holiday --
  Diwali) \\
\bottomrule
\end{tabular}
\caption{Timeline mapped to CS7.404 course milestones.}
\label{tab:timeline}
\end{table}

Assign 3 (due Oct 8), Quiz 2 (Oct 22) and Assign 4 (due Nov 2) overlap
this plan, which is why the real dataset freezes before the busiest
coursework stretch.

% ============================================================================
\section{Cut List and Open Items}
\label{sec:cutlist}
% ============================================================================

\textbf{Cut in this order if behind schedule:} soft matting
$\rightarrow$ Retinex $\rightarrow$ plastic stacked-layer cases
$\rightarrow$ bulk severity sweeps.

\textbf{Must survive:} real captures, DCP plus 2--3 classical
alternatives, matched synthetic data, and the ranking-agreement
comparison.

\textbf{To confirm with the TA / instructor:}
\begin{itemize}[nosep]
  \item Whether library calls (guided filter, CLAHE, Wiener, metrics)
        are acceptable for baselines.
  \item Whether the reframing from a DCP-robustness study to a
        multi-method, multi-medium comparison is acceptable, since the
        submitted proposal was DCP-centric.
  \item Exact Report 1 requirements.
\end{itemize}

\end{document}
